% !Mode:: TeX:UTF-8
\section{课题来源及研究的目的和意义}
\subsection{课题来源或研究背景}
近年来，随着智能交通系统（ITS）以及智能网联汽车（Connected and Automated Vehicles, CAV）技术的不断发展推广，车辆之间的通过车车（V2V）与车路（V2I）网联信息互联成为现实可能。高速公路因其全封闭、车速高、车流量大等基础物理特征，整体交通流对道路容量的微小变化极度敏感。一旦前方发生突发事故并导致部分车道封闭，即形成了具有时空动态约束特性的交通"瓶颈区"。此时，后方驶入该区域的车辆将面临必须执行强制换道（Mandatory Lane Change, MLC）的严峻挑战。在传统的纯人工驾驶（Human-Driven Vehicles, HDV）场景下，驾驶员的视距受限于前车遮挡与道路线形，难以提前获知事故确切位置，往往导致对前方异常状况的反应滞后。在接近瓶颈区并尝试汇入相邻车道时，人工驾驶车辆之间会由于路权争夺而产生激烈的相互博弈与冲突，从而打破原本稳定的交通流运行状态，引发"幽灵拥堵"、容量跌落甚至链式追尾等二次事故风险。

V2X 通信技术的引入，赋予了 CAV 超视距感知能力，使其能够在极大范围内提前捕获事故报警与位置信息，从而在长距离前置阶段即可从容进行换道准备。但是，在车流密度较高和市场渗透率动态变化的混合交通流环境下，智能体如何在确保安全的前提下，找到通行效率与操作舒适度的最佳平衡点，依然是一个高度非线性的博弈决策问题。

值得注意的是，在以往研究中，研究者常常将交通事故区域视为恒定不变的静态时空屏障。事故从发生到全局广播明显包含多个阶段——突发初期的局部信息受限期，与全局 V2X 广播激活后的有序疏导期——在这两个截然不同的阶段下，车辆的决策机制理应有着不同的自适应要求。因此，深耕事故路段下多阶段的协同演进机理，是实现自动驾驶技术规模化落地必须攻克的关键核心难点。

\subsection{研究的目的及意义}
本研究旨在构建一种基于多智能体混合博弈与 V2X 通信的高速公路事故路段 CAV 换道决策模型。目的在于突破传统阈值控制与单一博弈结构的局限，引入基于 8 维原子特征的混合博弈换道机制。

研究意义表现在：
1. 理论层面：细化了突发交通事故下的信息传播阶段，并将感知延迟、通信丢包、局部密度等真实环境因素融入车辆决策特征的计算中，提出混合博弈框架（静态 Nash / Stackelberg / Nash 合作）按后车异质性自适应选择博弈结构。同时，利用最大熵逆强化学习（MaxEnt IRL）从真实驾驶数据中学习决策权重，丰富了数据驱动博弈决策模型在微观交通流仿真中的应用。
2. 应用层面：模型可实现智能网联车辆在紧急情况下的自适应协同避让，在高达 3600 pcu/h 的高密度混合交通流中实现零碰撞，显著提高通行效率，为未来自动驾驶汽车在复杂高速场景中的规控系统设计提供参考依据，并通过消融实验量化 V2X 通信、博弈决策和 Python 控制各自对交通效率与安全的贡献。

\section{国内外在该方向的研究现状及分析}
\subsection{国外研究现状}
当前，国外研究者在智能汽车换道控制及自动驾驶决策调度领域已进行了大量而深入的探索。传统的微观交通流换道模型，如由 Gipps 提出的安全间距模型，以及后续基于间隙接受理论发展的 MOBIL 模型，主要聚焦于基于规则与本车绝对动力学表现的保守决策体系，较少关注与周边车辆交互时展现的主动协商策略。

因此，越来越多的国际权威学者开始应用博弈论来系统分析混合交通环境下的并发冲突与联合决策。Talebpour 等人提出了基于非合作与合作博弈交替计算的方法，以计算微观驾驶员的收益并精准预测换道意图。Yan 等人提出了多车辆博弈框架（IEEE Trans. Intelligent Vehicles, 2023），在混合交通中区分 CAV-HV（非合作 Stackelberg）和 CAV-CAV（部分合作博弈）。Huang 等人（IEEE Trans. ITS, 2024）提出了基于相对驾驶风格（RDS）的变结构博弈模型，根据驾驶风格的差异在静态博弈与 Stackelberg 博弈之间自适应切换，在 NGSIM 数据集上达到 98\% 的准确率。Wang 等人（IEEE Trans. ITS, 2026）提出了博弈论强化学习（GTRL）框架，对 CAV-HV 使用 Stackelberg 交互，对 CAV-CAV 使用 Nash 合作博弈。Liu 等人（IEEE Trans. VT, 2024）引入 Interaction Orientation 概念，提出了混合策略博弈替代传统纯策略，为行为不确定性建模提供了新思路。

\subsection{国内研究现状}
近年来，在国内"交通强国"战略及 5G-V2X 车路协同技术等新基建政策环境的大力扶持与驱动下，我国科研机构与高校在智能网联汽车协同控制方面同样进展迅速。

诸多基于人工智能交叉学科的实验室基于深度强化学习和多智能体协同（MADRL）框架，提出了大量能够适应高密度车流的自适应换道决策方法。He 等人（Scientific Reports, 2025）基于 K-Means 聚类将人类驾驶风格分为激进/普通/保守三类，并构建了多目标优化的换道模型，在 SUMO 平台上验证了在不同 CAV 渗透率下的安全性。Li 等人（IEEE Trans. IV, 2024）提出了结合博弈论与深度最大熵逆强化学习的车辆交互行为建模方法，从自然驾驶数据中学习博弈收益参数。Wu 等人（IEEE OJITS, 2024）将 Stackelberg 博弈与轨迹规划结合，实现了基于主从博弈的换道控制。不过，多数现有微观仿真体系倾向于将博弈结构固定化，对高度时变的通信质量退化与不同阶段自适应切换的研究尚不充分。

\subsection{国内外文献综述的简析}
通过系统梳理上述国内外研究轨迹可以看出：基于博弈论的车辆换道模型在理论严谨性上已得到学界的广泛认可。同时，V2X 技术的全面介入能显著改善微观通行能力并消除安全隐患。

然而，需要指出的是，当前文献仍存在以下研究空白：首先，严重匮乏对事故环境不同演进阶段（突发期与有序期）的微观切片研究，多数算法未能根据时序赋予不同的权重切换；其次，大量决策模型使用复合特征（如将速度比、紧迫度和车道压力加权为单一效率得分），导致内层系数与外层 PAYOFF_WEIGHTS 两层权重相互干扰，IRL 难以学出干净的物理解释。最后，在混合交通环境下，人类驾驶与 CAV 之间的交互博弈模式——CAV 应假设后车是理性的博弈参与方，还是应假设后车是不配合的 "非理性" 对象——现有文献大多采用单一博弈结构，缺乏按后车异质性自适应选择博弈结构的混合方案。

本课题针对上述空缺，提出包含 8 维原子特征空间 + 混合博弈结构 + 低渗透率自适应机制的完整方案，具备强烈的填补该方向空缺的学术价值。

\section{前期的理论研究与试验论证工作的结果}

\subsection{基于 8 维原子特征的博弈换道模型构建}
本课题在前期工作中，首先以 4 维复合特征模型为基础进行了仿真探索，发现 eff = 0.55·speed\_ratio + 0.45·urgency 式的复合特征导致 IRL 学出的权重难以解释。随后将特征空间重构为 8 维原子特征，每个维度具备独立物理语义，IRL 学出的权重可直接解释为"该维度的决策偏好"。

\textbf{特征空间定义}：对每个(本车动作, 后车反应)组合提取 8 个独立特征：

\begin{equation}
\mathbf{f}(a_{ego}, a_{fol}) = [speed, urgency, pressure, safe, coop, social, density, lc\_cost]
\end{equation}

各特征含义为：speed 表示目标车道速度收益；urgency 表示距事故紧迫度；pressure 表示当前车道受压程度；safe 表示综合 TTC 安全评分；coop 表示后车协同奖励；social 表示换道对后车 TTC 冲击；density 表示局部车流密度；lc\_cost 表示换道固有代价。关键在于，不换道时 urgency、pressure、social、density 直接置零，形成换道与不换道的非对称特征结构。

\textbf{混合博弈框架}：针对后车的异质性，设计了 3 种后车类型，各自对应不同的博弈结构：

\begin{itemize}
    \item Type 0（20\%）自私型：无协同动机，使用\textbf{静态 Nash} 博弈——双方同时决策，softmax 迭代收敛到混合策略均衡
    \item Type 1（60\%）合作型：基础协同意愿，使用\textbf{Stackelberg} 博弈——本车宣布动作后后车选择最优反应
    \item Type 2（20\%）高合作型：强协同意识，使用\textbf{Nash 合作} 博弈——最大化联合收益矩阵
\end{itemize}

\textbf{Social Impact 特征}：量化本车换道对目标车道后车 TTC 安全空间的侵占，social = max(0, (fol\_old\_ttc - fol\_new\_ttc) / fol\_old\_ttc)。当后车 TTC 因本车切入骤降时，正权重会降低换道收益，抑制强行插入——本研究将 Liu et al. (2024) 的 Interaction Orientation 概念首次量化为可计算的核特征并纳入博弈收益函数。

\subsection{混合交通仿真平台搭建}
为验证所提理论模型的有效性，依托开源微观交通流模拟引擎 Eclipse SUMO 构建了具备物理特性的交通数字孪生实验平台：

1. \textbf{场景架构}：仿真平台成功复现了 3.5km 三车道高速公路中间车道事故封堵的瓶颈场景。在通信链路中引入了丢包率模型与感知噪声/延迟，最大限度贴合实际工程环境。

2. \textbf{三种异质人类车}：为了实现真实的混合交通环境，基于 He et al. (2025) 的做法将人类驾驶员分为保守型（sigma=0.3, minGap=3.0m）、普通型（sigma=0.5, minGap=2.5m）和激进型（sigma=0.7, minGap=1.8m），克服了传统单一人型不能表现驾驶风格差异的缺陷。

3. \textbf{低渗透率自适应机制}：针对低 CAV 占比时博弈模式假设不成立的问题，设计了三种补偿机制——对抗性偏差（后车=人类时先验向不配合方向偏移）、互惠记忆（跟踪人类后车的合作历史，动态调整下次相遇的先验）和渗透率自适应阈值（min\_gain = base $\times$ (1.0 - 0.6 $\times$ (1 - CAV\_PENETRATION))），有效缓解了低渗透率下的决策失效。

\subsection{大规模仿真实验与结果分析}
课题组前期完成了 4 模型 $\times$ 4 密度 $\times$ 6 渗透率 = 96 组的完整消融实验（8 核并行）。四个模型构成三层消融：Game (Ours, 全功能) $\rightarrow$ No-V2X (关 V2X 通信) $\rightarrow$ Rule-Based (关博弈, 多条件规则 TTC 前向 1.5s + 后向 1.0s) $\rightarrow$ SUMO Default (关 Python 控制, SL2015 原生)。

\textbf{Game (Ours) 渗透率扫描}（通过车辆数 / 换道次数）：

\begin{table}[H]
\centering
\caption{Game (Ours) 完整渗透率扫描}
\label{tab:penetration}
\small
\begin{tabular}{c|cccc}
\hline
渗透率 & 1200 pcu/h & 2000 pcu/h & 2800 pcu/h & 3600 pcu/h \\
\hline
0\% & 77 / 0 & 109 / 0 & 97 / 0 & 61 / 0 \\
10\% & 77 / 4 & 118 / 6 & 172 / 18 & 214 / 22 \\
30\% & 79 / 8 & 123 / 28 & 159 / 33 & 213 / 37 \\
50\% & 74 / 23 & 122 / 35 & 171 / 61 & 211 / 74 \\
70\% & 78 / 33 & 124 / 58 & 167 / 99 & 214 / 134 \\
100\% & 77 / 43 & 126 / 82 & 173 / 115 & 204 / 205 \\
\hline
\end{tabular}
\end{table}

\textbf{全模型对比 @ 100\% CAV}：

\begin{table}[H]
\centering
\caption{全模型对比 @ 100\% CAV}
\label{tab:model_comparison}
\small
\begin{tabular}{c|cccc|c}
\hline
模型 & 1200 & 2000 & 2800 & 3600 & 碰撞 \\
\hline
Game (Ours) & 77 & \textbf{126} & \textbf{173} & \textbf{204} & \textbf{0} \\
No-V2X & 77 & 126 & 173 & 203 & 2 \\
Rule-Based & 48 & 80 & 106 & 135 & 0 \\
SUMO Default & 61 & 35 & 41 & 57 & 2 \\
\hline
\end{tabular}
\end{table}

\textbf{消融层分析结论}：

\begin{itemize}
    \item L1（Game vs No-V2X）：通过量一致，但 No-V2X 有碰撞——V2X 通信的贡献在\textbf{安全}层面
    \item L2（No-V2X vs Rule-Based）：No-V2X 在 2000 pcu/h 高出 57\%（126 vs 80）——博弈决策大幅提升效率
    \item L3（Rule-Based vs SUMO Default）：Rule-Based 在 2000 pcu/h 高出 128\%（80 vs 35）——Python 控制的价值
    \item 渗透率效应：从 10\% 至 100\% 通过量保持稳定（77-79 辆/1200 pcu/h），30\% CAV 已可获得全部收益
\end{itemize}

\subsection{MaxEnt IRL 数据驱动权重学习}
从 AD4CHE 数据集（大疆无人机采集的中国高速公路数据，360 万帧，3908 个换道片段）中学习 8 维权重。IRL 训练 30 轮，loss 从 4.488 收敛到 4.293。最终权重为：

\begin{equation}
\mathbf{w}_{IRL} = [0.15, 0.19, 0.18, 0.25, 0.02, 0.00, 0.11, 0.70]
\end{equation}

lc\_cost=0.70 为权重主导——真人司机天然不爱频繁换道；safe=0.25 为第二高权重——安全始终普遍考虑。IRL 权重在全场景实验中与手工权重效果相当，印证了数据驱动方案的可替代性。

\section{学位论文的主要研究内容、实施方案及其可行性论证}
\subsection{主要研究内容}
\subsubsection{混合博弈框架的数学建模与可学习特征空间设计}
进一步完善 8 维原子特征空间的理论推导，建立从感知输入到特征输出到博弈收益的完整映射链。引入自适应权重机制——权重向量 $\mathbf{w}$ 随局部密度 $\rho$ 和紧迫度 $U_{urgency}$ 动态调整：

\begin{equation}
\mathbf{w}_{safe} = \mathbf{w}_{safe}^{base} + 0.10 \cdot \rho + 0.05 \cdot U_{urgency}
\end{equation}

\begin{equation}
\mathbf{w}_{speed} = \max(0.05, \mathbf{w}_{speed}^{base} - 0.10 \cdot U_{urgency})
\end{equation}

同时深化混合博弈求解器的数学架构，推导三种博弈结构下（Nash 均衡、Stackelberg 均衡、Nash 合作解）的均衡存在性与收敛条件。

\subsubsection{多层递进消融实验框架}
基于已运行的 96 组实验，进一步扩展消融维度。考虑增加：（1）心理因素——后车行为不仅受物理状态影响，还受"认知层级"（Level-k）影响；（2）针对不同围堵密度区间的局部拥堵演化特性做精细化分析，建立从宏观交通流密度到微观博弈参数的连续映射。

\subsubsection{V2X 通信质量退化的鲁棒性边界研究}
在不完美通信信道模型的基础上进一步扩展，将 V2X 丢包率作为实验变量纳入实验矩阵（5\% - 20\%），测量通信质量退化对不同渗透率下系统整体安全性的影响，寻找博弈模型保持鲁棒性的 V2X 通信质量边界。

\subsection{实施方案及其可行性论证}
实施步骤如下：
\begin{enumerate}
    \item \textbf{理论深化}：针对三种博弈结构的均衡等性质给出推导和边界条件证明；对在线 TD 学习的收敛性作出分析
    \item \textbf{实验扩展}：在现有 4 模型 $\times$ 4 密度 $\times$ 6 渗透率实验的基础上，增加通信丢包率维度和 Level-k 认知维度，进一步验证模型的鲁棒性与泛化能力
    \item \textbf{数据挖掘}：对前期输出的大量仿真轨迹数据进行深度挖掘，提取队列形成与消散的时空特征、换道冲突热点区域、延误分布等宏观指标
    \item \textbf{论文撰写}：统筹分析结果和可视化图表，完学位论文的撰写工作
\end{enumerate}

前期已搭建完成全部的仿真代码框架和实验运行脚本，具备快速扩展实验的能力。SUMO + TraCI + Python 的技术路线已被大量国际研究验证，在高性能计算集群的支撑下，论文核心实验的补充工作可在一个月内全部完成。

\section{论文进度安排，预期达到的目标}
\subsection{进度安排}
\begin{itemize}
    \item \textbf{第一季度（模型设计与核心代码重构）：}
    完成 8 维特征空间的理论公式推导，完善混合博弈求解器的数学框架。在现有代码的基础上加入 Level-k 认知层级和 V2X 通信质量退化变量。
    \item \textbf{第二季度（大规模对比实验运行）：}
    运行扩展后的实验矩阵（含通信质量维度和认知维度），完成全部仿真数据采集。利用高分辨率浮动车轨迹数据提取多项宏微观评价指标。
    \item \textbf{第三季度（数据分析和论文初稿）：}
    完成实验数据的统计分析、可视化图形制作、消融结果深度解读。完成核心章节的学术逻辑撰写。
    \item \textbf{第四季度（论文修订和答辩准备）：}
    完成学位论文初稿，经过修回、校对进而定稿。按校方时间节点准备有关中期检查及答辩的相关材料。
\end{itemize}

\subsection{预期达到的目标}
\textbf{主要预期目标：}
1. 建立并标定一套涵盖混合博弈结构、8 维原子特征、V2X 信道质量波动的混合交通换道控制模型。该模型在安全性与效率性指标上显著优于 SUMO Default 基线系统。
2. 通过消融实验证实博弈控制算法在最高密度（3600 pcu/h）场景下保持零碰撞，同时将通行效率提升 150\% 以上（2000pcu/h 时 Game 126 vs SUMO Default 35 辆）。
3. 利用 IRL 从真实驾驶数据中学习决策权重，实现数据驱动与仿真实验的闭环。

\textbf{成果形式：}
撰写完成一篇高质量、逻辑缜密的硕士学位论文。力争在 IEEE ITSC、TRB 年会或国内核心期刊上投稿 1-2 篇科研论文。

\section{学位论文预期创新点}
1. \textbf{提出面向异质性后车的混合博弈结构}：基于后车驾驶风格（自私/合作/高合作）自适应选择博弈结构（静态 Nash / Stackelberg / Nash 合作），克服了传统单一博弈结构不能适应不同后车类型的缺陷。

2. \textbf{设计 8 维原子特征空间}：打破复合特征导致的两层权重耦合问题，每个特征维度具备独立物理语义。IRL 学出的权重可直接解释为决策偏好，为数据驱动的可解释博弈决策提供了理论基础。

3. \textbf{Social Impact 特征}：首次将换道对后车 TTC 的空间侵占量化为可计算特征并纳入博弈收益函数，使换道决策不仅考虑"我能不能更快"，还考虑"我换了之后后车会不会被迫急刹"。

4. \textbf{低渗透率自适应三机制}：针对低 CAV 占比时博弈模式失效的问题，提出对抗性偏差、互惠记忆和渗透率自适应阈值三套协同机制。

5. \textbf{逐层消融实验框架}：通过关闭 V2X→博弈→Python 控制，逐层量化各功能模块的贡献，为混合交通 CAV 系统设计提供定量设计参考。

\section{为完成课题已具备和所需的条件、外协计划及经费}
\textbf{已具备条件：}本课题已成功搭建 SUMO-TraCI-Python 联合仿真平台，并基于该平台完成了 4 模型 $\times$ 4 密度 $\times$ 6 渗透率 = 96 组的全量消融实验。所有仿真代码（~3000 行）已经开发部署完毕并在 GitHub 开源。前期实验数据包含通过量、换道次数、碰撞计数、队列长度、行程时间等关键指标，具备继续深化研究的数据底座。

\textbf{所需条件及经费：}对于更高阶的扩展实验（通信质量退化维度），需要在服务器租赁部分增加少量算力预算。该部分由课题总经费拨发。不涉及超出常规的外部协作计划。

\section{预计研究过程中可能遇到的困难、问题，以及解决的途径}
\textbf{潜在困难与问题：}目前模型在每一时间步对每辆 CAV 执行完整的三种博弈结构求解，在高密度场景下计算开销较大。在饱和度极高（360辆全 CAV）时，单轮仿真耗时可达 10 分钟以上。

\textbf{解决途径：}采用决策减频优化（每 3 步执行一次博弈，已实现）和速度缓存（同类型车辆最大速度只查一次 TraCI，已实现）；后续可引入离线策略预训练（如利用强化学习替代在线博弈求解），将在线运算量进一步消减。

\nocite{*}
\makeatletter
\renewcommand{\bibsection}{}
\makeatother
\section{主要参考文献}
\bibliographystyle{heuthesis}
\bibliography{reference}
